\documentclass[letterpaper,10pt]{article}

\usepackage{latexsym}
\usepackage[empty]{fullpage}
\usepackage{titlesec}
\usepackage{marvosym}
\usepackage[usenames,dvipsnames]{color}
\usepackage{verbatim}
\usepackage{enumitem}
\usepackage[hidelinks]{hyperref}
\usepackage{fancyhdr}
\usepackage[brazilian]{babel}
\usepackage{tabularx}
\input{glyphtounicode}

\usepackage[sfdefault]{roboto}

\pagestyle{fancy}
\fancyhf{}
\fancyfoot{}
\renewcommand{\headrulewidth}{0pt}
\renewcommand{\footrulewidth}{0pt}
\setlength{\footskip}{5pt}

\addtolength{\oddsidemargin}{-0.5in}
\addtolength{\evensidemargin}{-0.5in}
\addtolength{\textwidth}{1in}
\addtolength{\topmargin}{-.5in}
\addtolength{\textheight}{1.0in}

\urlstyle{same}
\raggedbottom
\setlength{\tabcolsep}{0in}

% Configuração para justificar o texto em blocos e sem recuo
\setlength{\parindent}{0pt}
\setlength{\parskip}{7pt}

\titleformat{\section}{\Large\bfseries\scshape\raggedright}{}{0em}{}[\titlerule]

\pdfgentounicode=1

\newcommand{\resumeItem}[1]{\item\small{#1}}
\newcommand{\resumeSubheading}[4]{
\vspace{-1pt}\item
  \begin{tabular*}{0.97\textwidth}[t]{l@{\extracolsep{\fill}}r}
    \textbf{#1} & #2 \\
    \textit{#3} & \textit{#4} \\
  \end{tabular*}\vspace{-7pt}
}
\renewcommand\labelitemii{$\vcenter{\hbox{\tiny$\bullet$}}$}
\newcommand{\resumeSubHeadingList}{\begin{itemize}[leftmargin=0.15in, label={}]}
\newcommand{\resumeSubHeadingListEnd}{\end{itemize}}

\begin{document}

\begin{center}
  \textbf{\Huge Gabriel Frigo} \\ \vspace{4pt}
  \small +55 11 95570-2000 $|$ \href{mailto:gabriel.frigo4@gmail.com}{gabriel.frigo4@gmail.com} $|$
  \href{https://linkedin.com/in/gabriel-frigo-b6727b275/}{linkedin.com/in/gabriel-frigo} $|$
  \href{https://github.com/GabrielFrigo4}{github.com/GabrielFrigo4} $|$
  \href{https://gabrielfrigo.dev.br}{gabrielfrigo.dev.br}
\end{center}

\vspace{6pt}

\textbf{Assunto: Carta de Apresentação e Candidatura -- Engenharia de Software}

Prezada Equipe de Engenharia e Recrutamento,

Escrevo para expressar meu interesse em oportunidades de atuação em \textbf{Engenharia de Software, Sistemas e Infraestrutura}. Sou estudante de \textbf{Ciência da Computação na UFABC}, com um perfil que une a densidade matemática e acadêmica à engenharia de software pragmática, sistemas de baixo nível e infraestrutura em nuvem.

Minha abordagem técnica transita com naturalidade entre \textbf{a teoria acadêmica e a produção do mundo real}: encontro beleza tanto no rigor da \textbf{Matemática e Pesquisa Operacional} (onde aplico Teoria dos Grafos, Algoritmos de Fluxo e Combinatória) quanto na engenharia aplicada moderna (IaC com OpenTofu, microsserviços em Go, Docker/Podman e nuvem). Da mesma forma, domino o ecossistema corporativo contemporâneo sem abrir mão do apreço pelos fundamentos artesanais de sistemas — criando de scripts shell e instalações nativas em \textbf{FreeBSD com Jails (Bastille)} a códigos de alto desempenho em \textbf{C/C++, Lua, Got e padrões POSIX}.

No lado prático, desenvolvi o \textbf{OptiLaser} — motor de decisão logística e pesquisa operacional (VRPTW) que resolve problemas complexos de otimização combinatória em $<3$s em \textbf{Go e CGo} com \textbf{Google OR-Tools}, multiplataforma (\textbf{FreeBSD/Linux}), banco \textbf{PocketBase} embarcado e provisionamento IaC via \textbf{OpenTofu} na \textbf{Magalu Cloud} sobre \textbf{Podman com Rocky Linux}. Também implementei do zero um \textbf{servidor HTTP concorrente em C} sobre GNU/BSD Sockets e domino linguagens como \textbf{C, C++, Rust e Assembly}, que aplico na instrução técnica no \textbf{Green Team Hacker Club} e em robótica competitiva.

Minha base analítica foi forjada por participações de destaque em competições, como a \textbf{Medalha de Ouro (13ª posição nacional) na OBMEP}, a \textbf{Maratona SBC de Programação} e o \textbf{ICPC}, além de pesquisas de Iniciação Científica em \textbf{Fluxos em Redes (CNPq)} e \textbf{Teoria de Ramsey (PICME)}. Aliando um forte perfil \textbf{autodidata} à cultura \textbf{AI-First}, utilizo a Inteligência Artificial não apenas para acelerar a entrega de código, mas como um interlocutor para o aprendizado contínuo através do \textbf{método socrático} — questionando premissas, explorando \textit{trade-offs} de arquitetura e aprofundando conceitos técnicos com velocidade e rigor.

Acredito que essa versatilidade — unindo modelagem matemática, sistemas de baixo nível, infraestrutura escalável e aprendizado acelerado — me qualifica para gerar impacto imediato em desafios técnicos de alta complexidade. Estou à disposição para uma conversa.

Atenciosamente,

\textbf{Gabriel Frigo}

\end{document}
